%%% FIELD proof-unbalanced-frontier
Write \(C=(h,g)\) and \(C'=(h',g')\), and adopt the convention that \(\binom{x}{y}=0\) for integers \(x,y\) unless \(0\leq y\leq x\).  The causal quantity \(\tau_R\) is fixed, rather than random, because \(R\) is fixed; membership in \(\mathcal R_0(n,K)\), by Definition~\(\mathrm{def:legal\mbox{-}response\mbox{-}tables}\) and Assumption~\(\mathrm{ass:weak\mbox{-}null}\), gives \(\tau_R=0\).  The kernel \(P_R^{(n_1)}\) below is instead an observed-profile probability induced solely by complete randomization.  We first verify that kernel and then characterize its common level sets.

For fixed \(R\), a treated-count matrix \(T\) records that exactly \(T_{ab}\) of the \(R_{ab}\) labeled units of type \((a,b)\) are treated.  There are
\[
 \prod_{a,b\in\mathcal A_K}\binom{R_{ab}}{T_{ab}}
\]
assignments giving that matrix.  The row constraints on \(T\) are exactly the requirement that the treated histogram be \(h\), and the column constraints on \(R-T\) are exactly the requirement that the control histogram be \(g\).  Distinct matrices describe disjoint assignments.  Hence summing these products over \(\mathcal T_R^{(n_1)}(h,g)\) gives \(W_R^{(n_1)}(h,g)\).  Complete randomization is uniform on the \(\binom n{n_1}\) treated sets; this denominator is positive because \(1\leq n_1\leq n-1\).  Consequently
\[
 P_R^{(n_1)}(h,g)=\binom n{n_1}^{-1}W_R^{(n_1)}(h,g).
 \tag{1}
\]

Let \(r=h+g\).  Then \(r\in\mathbb Z_+^K\) and \(\sum_a r_a=n_1+n_0=n\).  The diagonal table \(D_r\) has mass \(n\) by Assumption~\(\mathrm{ass:table\mbox{-}mass}\), and all of its units have zero individual effect, so Assumption~\(\mathrm{ass:weak\mbox{-}null}\) gives \(D_r\in\mathcal R_0(n,K)\).  Directly from the treated-count constraints and Definition~\(\mathrm{def:assignment\mbox{-}count}\),
\[
 W_{D_r}^{(n_1)}(u,v)
 =\mathbf 1\{u+v=r\}\prod_{a\in\mathcal A_K}\binom{r_a}{u_a}.
 \tag{2}
\]
In particular,
\[
 W_{D_r}^{(n_1)}(h,g)=Q_r(h)>0,
 \tag{3}
\]
because \(0\leq h_a\leq r_a\) for every \(a\).  If \(C\sim_{n_1,n_0}C'\), equality of the probabilities at \(D_r\), together with (1)--(3), first forces
\[
 h'+g'=r
 \tag{4}
\]
and then forces
\[
 Q_r(h)=Q_r(h').
 \tag{5}
\]

We next show that legal response tables having two heterogeneous units recover the arm-imbalance outer product.  Fix distinct \(a,b\in\mathcal A_K\) with \(r_a,r_b>0\), and define
\[
 R^{ab}=D_{r-e_a-e_b}+E_{ab}+E_{ba}.
 \tag{6}
\]
The diagonal filler in (6) is nonnegative.  The table has mass \(n\), and the effects of its two heterogeneous units sum to \((a-b)+(b-a)=0\).  Thus (6) belongs to \(\mathcal R_0(n,K)\), and Definition~\(\mathrm{def:heterogeneous\mbox{-}count}\) gives \(q(R^{ab})=2\).  The two special units contribute pooled observed histogram \(e_a+e_b\) if they are both treated or both controlled, \(2e_a\) if the \((a,b)\) unit alone is treated, and \(2e_b\) under the other mixed allocation.  Since the diagonal filler has pooled histogram \(r-e_a-e_b\), precisely the both-treated and both-controlled cases can realize a profile whose pooled histogram is \(r\).  Using
\[
 \binom{s-1}{t-1}=\frac{t}{s}\binom{s}{t},
 \qquad
 \binom{s-1}{t}=\frac{s-t}{s}\binom{s}{t}
 \tag{7}
\]
with the positive denominators \(s=r_a,r_b\), the two cases give
\[
 W_{R^{ab}}^{(n_1)}(h,g)
 =Q_r(h)\frac{h_ah_b+g_ag_b}{r_ar_b}.
 \tag{8}
\]
Define the integer imbalance vectors
\[
 x=h-g=2h-r,
 \qquad
 x'=h'-g'=2h'-r,
 \tag{9}
\]
where (4) is used for the second equality defining \(x'\).  Since
\[
 h_ah_b+g_ag_b=\frac{r_ar_b+x_ax_b}{2},
 \tag{10}
\]
equivalence, (1), (5), and (8)--(10) imply
\[
 x_ax_b=x'_ax'_b
 \quad\text{for distinct }a,b\text{ with }r_a,r_b>0.
 \tag{11}
\]

The diagonal entries of the same outer product are obtained by a second placement of the symmetric switch.  Fix \(a\) with \(r_a\geq2\), choose any \(b\neq a\), which is possible because \(K\geq2\), and put
\[
 \widetilde R^{a\mid b}=D_{r-2e_a}+E_{ab}+E_{ba}.
 \tag{12}
\]
As in (6), this is a nonnegative table of mass \(n\), its two heterogeneous effects cancel, and \(q(\widetilde R^{a\mid b})=2\).  Exactly one treatment pattern of the two special units supplies the missing pooled histogram \(2e_a\): the \((a,b)\) unit is treated and the \((b,a)\) unit is controlled.  Therefore the coefficient formula gives
\[
 \begin{aligned}
 W_{\widetilde R^{a\mid b}}^{(n_1)}(h,g)
 &=\prod_c\binom{r_c-2\mathbf 1\{c=a\}}
                         {h_c-\mathbf 1\{c=a\}}\\
 &=Q_r(h)\frac{h_ag_a}{r_a(r_a-1)}.
 \end{aligned}
 \tag{13}
\]
The division in (13) is legitimate because \(r_a\geq2\); the equality also covers the boundary cases \(h_a=0\) and \(h_a=r_a\) under the binomial convention.  Equivalence, (1), (5), and (13) imply \(h_ag_a=h'_ag'_a\), and hence
\[
 x_a^2=r_a^2-4h_ag_a=r_a^2-4h'_ag'_a=(x'_a)^2.
 \tag{14}
\]
If \(r_a=0\), both imbalance coordinates are zero.  If \(r_a=1\), nonnegativity and integrality with \(h_a+g_a=h'_a+g'_a=1\) imply \(x_a^2=(x'_a)^2=1\).  Thus (14) holds for every \(a\).  Products involving a coordinate with \(r_a=0\) vanish, so (11) and (14) together yield
\[
 xx^{\mathsf T}=x'(x')^{\mathsf T}.
 \tag{15}
\]

The unequal allocation now eliminates the sign ambiguity in (15).  The profile constraints give
\[
 \sum_a x_a=\sum_a x'_a=n_1-n_0=:\delta,
 \qquad \delta\neq0.
 \tag{16}
\]
Thus \(x\neq0\).  Choose \(j\) with \(x_j\neq0\).  The \((j,j)\) entry of (15) gives \(x'_j=\varepsilon x_j\) for some \(\varepsilon\in\{-1,1\}\); the \((j,a)\) entries then give \(x'_a=\varepsilon x_a\) for every \(a\).  Summing and using (16) gives \(\delta=\varepsilon\delta\), so \(\varepsilon=1\).  Therefore \(x'=x\), and (4) and (9) give
\[
 h'=\frac{r+x'}2=\frac{r+x}2=h,
 \qquad
 g'=\frac{r-x'}2=\frac{r-x}2=g.
 \tag{17}
\]
This proves that every \(\sim_{n_1,n_0}\)-class is a singleton; the converse implication \(C=C'\Rightarrow C\sim_{n_1,n_0}C'\) is immediate from the definition.

The preceding argument is also a constructive separator algorithm.  For distinct \(C,C'\), first use \(D_{h+g}\).  Equation (2) separates them if their pooled histograms differ, and (2)--(3) separate them if the pooled histograms agree but their diagonal coefficients differ.  Suppose instead that their common pooled histogram is \(r\) and (5) holds.  If some product in (11) differs, the explicitly defined table (6) separates the profiles by (8).  If all such off-diagonal products agree but the profiles remain distinct, then some square in (14) must differ: otherwise (15)--(17), whose only additional premise is the common nonzero sum \(\delta\), would imply \(C=C'\).  A coordinate with unequal squares must satisfy \(r_a\geq2\), since the square is forced to be zero for \(r_a=0\) and one for \(r_a=1\); the table (12) then separates the profiles by (13).  Each table in this algorithm is legal and has respectively zero or two heterogeneous units.  It follows that
\[
 s_\star(n_1,n_0,K)\leq2.
 \tag{18}
\]

If \(f\in G_\star^{(n_1,n_0)}\), its defining probability identity implies \(C\sim_{n_1,n_0}f(C)\) for every \(C\).  The singleton result forces \(f(C)=C\) for every profile, and hence
\[
 G_\star^{(n_1,n_0)}=\{\operatorname{id}\}.
 \tag{19}
\]

It remains to prove that (18) is sharp.  The assumptions permit no unequal positive allocation when \(m=1\), so every admissible instance has \(m\geq2\).  Define
\[
 \ell=\max\{0,n_1-m\},
 \qquad
 u=\min\{m,n_1\}.
 \tag{20}
\]
If \(n_1<m\), then \((\ell,u)=(0,n_1)\); if \(n_1>m\), then \((\ell,u)=(n_1-m,m)\).  Using \(1\leq n_1\leq2m-1\) and \(n_1\neq m\), both cases give
\[
 0\leq\ell<u\leq m,
 \qquad \ell+u=n_1.
 \tag{21}
\]
Because \(K\geq2\), the levels \(0,1\) belong to \(\mathcal A_K\).  Define
\[
 \begin{aligned}
 h&=\ell e_0+u e_1,
 &g&=(m-\ell)e_0+(m-u)e_1,\\
 h'&=u e_0+\ell e_1,
 &g'&=(m-u)e_0+(m-\ell)e_1.
 \end{aligned}
 \tag{22}
\]
Equations (21)--(22) show that the treated totals are \(n_1\), the control totals are \(2m-n_1=n_0\), all entries are nonnegative, and the two profiles are distinct.  Both have pooled histogram
\[
 r=me_0+me_1.
 \tag{23}
\]
If a legal table \(R\) has \(q(R)=0\), nonnegativity in Definition~\(\mathrm{def:heterogeneous\mbox{-}count}\) forces every off-diagonal entry to vanish, so \(R\) is diagonal.  By (2), both counts in (22) vanish unless \(R=D_r\); for that table,
\[
 W_{D_r}^{(n_1)}(h,g)
 =\binom m\ell\binom mu
 =\binom mu\binom m\ell
 =W_{D_r}^{(n_1)}(h',g').
 \tag{24}
\]
Thus no legal table with zero heterogeneous units separates this pair.  There is no legal table with exactly one heterogeneous unit: all diagonal units contribute zero to the weak-null sum, while the sole off-diagonal unit of type \((a,b)\) contributes \(a-b\neq0\), contradicting Assumption~\(\mathrm{ass:weak\mbox{-}null}\).  Therefore the pair (22) has separator number at least two.  Combining this lower bound with (18) proves
\[
 s_\star(n_1,n_0,K)=2.
 \tag{25}
\]

As an exact smallest-instance check, take \((n,K,n_1,n_0)=(4,2,1,3)\).  Then (22) gives \((h,g)=((0,1),(2,1))\) and \((h',g')=((1,0),(1,2))\).  Under \(D_{2e_0+2e_1}\), both assignment counts equal \(2\), and under every other diagonal table both vanish.  The legal table
\[
 2E_{11}+E_{01}+E_{10}
\]
has two heterogeneous units and, by direct substitution in (13) with \(a=0\) and \(b=1\), gives the respective counts \(0\) and \(1\).  This verifies at the boundary both the explicit upper construction and the obstruction to a separator with fewer than two heterogeneous units.
